\documentclass{report}
\usepackage{amsmath,amssymb}
\setlength{\parindent}{0mm}
\setlength{\parskip}{1em}
\begin{document}
\begin{center}
\rule{6in}{1pt} \
{\large Veselin Dobrev \\
{\bf Iterative Relaxation of High-order Curvilinear Meshes}}

Center for Applied Scientific Computing \\ Lawrence Livermore National Laboratory \\ P O Box 808 \\ L-560 \\ Livermore \\ CA 94551
\\
%{\tt tzanio@llnl.gov}\\
{\tt dobrev1@llnl.gov}\\
Robert Anderson\\
%Veselin Dobrev\\
Tzanio Kolev\\
Robert Rieben\end{center}

In this talk we discuss our experience with iterative methods for
relaxation of high-order curvilinear meshes, which aim to improve
the quality of the curved mesh elements through successive
iterates for the coordinates of the high-order mesh
nodes (control points). The motivation for our work is the need
for new mesh optimization algorithms in the development of
high-order finite element discretizations of shock hydrodynamics
problems in the Arbitrary Lagrangian-Eulerian (ALE) framework
[1], but the need for mesh improvement is present also in other
areas of computational science, and is of general interest. The
goal of mesh relaxation is to define a new computational mesh
that is better than the current mesh with respect to some metric
related to mesh distortion. We will review several such metrics,
obtained by defining an appropriate high-order topological ``mesh
Laplacian'' operator. We will focus on harmonic-type mesh
relaxation, where a simple iterative scheme is applied to a
linear system involving the mesh Laplacian and the high-order
node coordinates. We will also investigate different
smoothers/preconditioners for this iteration, such as the
recently proposed polynomial method [2], as well as the effect of
different choices for the high-order basis. The harmonic approach
will be compared with other local and global high-order
optimization methods available in the Mesquite mesh quality
improvement toolkit [3].

[1] BLAST: High-order curvilinear finite elements for shock
hydrodynamics}, http://www.llnl.gov/CASC/blast.

[2] M.Berndt and N.Carlson, "Using Polynomial Filtering for Rezoning in
ALE", Talk at the International Conference on Numerical Methods for
Multi-Material Fluid Flows, Archachon, France, Sep 5--9, 2011, LA-UR
11-05015.

[3] Mesquite: Mesh Quality Improvement Toolkit,
http://www.cs.sandia.gov/optimization/knupp/Mesquite.html.


\end{document}
